% ---------------------------------------------------------------------
%  How to read this manual  (unnumbered front-matter chapter)
% ---------------------------------------------------------------------
\chapter*{How to Read This Manual}
\addcontentsline{toc}{chapter}{How to Read This Manual}
\markboth{How to Read This Manual}{How to Read This Manual}

\section*{Who this is for}

This manual documents the internals of \Daedalus{}, the automated
\msrjd{} Feynman-diagram engine.  It is written for someone who
understands the \emph{physics} --- stochastic differential equations, the
response-field (\msrjd) formalism, Feynman diagrams, and the loop expansion ---
but who has \emph{not} necessarily worked with the specific software libraries
the engine is built on.  The two audiences we keep in mind are:

\begin{itemize}[nosep]
  \item a new lab member who wants to \emph{use} the pipeline to compute
        cumulants for their own model, and
  \item a developer who needs to \emph{modify or extend} a specific stage and
        must understand exactly what the code does and why.
\end{itemize}

Wherever a piece of external tooling appears --- SageMath, the \code{nauty}
graph generator, \code{sympy}, \code{numba}, \code{scipy} --- we explain what it
is and how the engine uses it, from the ground up.  When we must choose between
saying too much and too little, we say too much.

\section*{How the manual is organised}

The book is divided into six parts that follow the data as it flows through the
pipeline, plus appendices:

\begin{description}[style=nextline,leftmargin=1.2em]
  \item[Part I --- Orientation] A bird's-eye view of the whole pipeline and a
        worked ``hello world'': from a text description of a model to a plotted
        cumulant in a dozen lines.
  \item[Part II --- Specifying a Model] How you describe a stochastic (S)PDE to
        the engine: the \code{TheoryBuilder} authoring API, the on-disk
        \file{.theory.py} file format, and the automatic rewriting of
        colored noise into an auxiliary Markovian field.
  \item[Part III --- Building the Perturbation Theory] The symbolic core:
        expanding the action into vertices, constructing the propagator,
        solving for the mean-field saddle, enumerating diagrams, assigning
        edge types, computing symmetry factors, and caching all of it.
  \item[Part IV --- Evaluating Diagrams: Temporal Models] The time-domain loop
        integrator (``Phase~J''), both the per-diagram and the grouped analytic
        mode-sum paths.
  \item[Part V --- Evaluating Diagrams: Spatial Models] The momentum-space
        integrator for stochastic PDEs: Symanzik polynomials, causal chambers,
        the heat-kernel inverse Fourier transform, and coupled multi-field
        dressing.
  \item[Part VI --- Assembly and the User Surface] How the per-diagram numbers
        are assembled into cumulants and moments, the user-facing
        \code{daedalus} API, the validation simulators, and the graphical
        theory-builder UI.
  \item[Appendices] A glossary, a consolidated primer on the external tools, a
        file-by-file index of the source tree, and a register of known issues
        and open questions.
\end{description}

You do not have to read linearly.  If you only want to \emph{use} the pipeline,
Part~I and the \code{daedalus} API chapter are enough.  If you are debugging a
specific stage, jump straight to its chapter; each is written to stand on its
own, with back-references where context is needed.

\section*{Typographic conventions}

\begin{itemize}[nosep]
  \item \code{Inline monospace} denotes code: a function, variable, keyword
        argument, or short snippet.
  \item \file{src/path/to/file.py} (navy monospace) denotes a path in the
        repository, always relative to the repository root.
  \item The \emph{first} appearance of a technical \term{term} is set in bold
        and defined nearby; every such term also appears in the glossary
        (Appendix~A).
  \item Multi-line code is shown in framed listings.  Python listings use
        syntax highlighting; console sessions and output are shown verbatim.
\end{itemize}

Three kinds of callout box punctuate the text:

\begin{note}
A clarification or a useful aside --- safe to skim, helpful to read.
\end{note}

\begin{gotcha}
A sharp edge: a non-obvious requirement, a platform-specific hazard, or a place
where the obvious thing is wrong.  Read these.
\end{gotcha}

\begin{defn}
A precise definition of a term or object the surrounding text depends on.
\end{defn}

\section*{Running the examples}

Every runnable example assumes the engine is imported through its convenience
layer:

\begin{lstlisting}
import daedalus as dd                 # notebooks/daedalus.py
model, mod = dd.load_theory('ou_quartic')
cfg = dd.Config(k=2, max_ell=2,
                external_fields=[('dx', 1), ('dx', 1)],
                parameters={'mu': 1.0, 'eps': 0.02, 'D': 1.0})
res = dd.run(model, cfg, mod)
\end{lstlisting}

\begin{gotcha}
The engine runs inside \textbf{SageMath}, not a plain CPython interpreter.  All
scripts must be launched with \code{sage -python yourscript.py}, and notebooks
must use the \emph{SageMath} Jupyter kernel.  Sage provides the exact symbolic
algebra and the multivariate polynomial rings the engine relies on; a stock
\code{python3} will not have them.  See Appendix~B for what Sage is and why it is
required.
\end{gotcha}

\section*{Getting the engine running}

Everything in this manual assumes a working \textbf{SageMath} with the project's
dependencies importable. If you are starting from nothing:
\begin{enumerate}
  \item \textbf{Install SageMath} (10.x). On macOS the binary app or a conda
        install both work; the engine only needs \code{import sage.all} to
        succeed.
  \item \textbf{Make the repository importable.} The example notebooks do this
        automatically by walking up to the directory that contains the
        \file{pipeline/} package and adding it (and \file{notebooks/}) to
        \code{sys.path}.
  \item \textbf{Use the project's environment and launcher.} The pipeline is
        developed in a dedicated conda environment (named \code{MSRJD\_diagrams}
        in this project) layered on Sage, and \emph{every} script is run with
        \code{sage -python}, never a bare \code{python3} (Appendix~B explains
        why).
  \item \textbf{Smoke-test} in under a second from the repository root:
\begin{lstlisting}[style=console]
$ sage -python -c "import sys; sys.path.insert(0,'notebooks'); import daedalus as dd; print(dd.list_theories()[:3])"
\end{lstlisting}
        If that prints a few theory names, the environment is green and you can
        go straight to the Quickstart (Chapter~\ref{ch:quickstart}).
\end{enumerate}

\begin{gotcha}
A stock miniforge/conda \emph{base} environment can ship a \textbf{broken NumPy}
(a \code{libgfortran} rpath problem) that makes \code{import numpy} fail under
Sage. If you hit an opaque NumPy import error, you are almost certainly in the
wrong environment --- activate the project's \code{MSRJD\_diagrams} env and try
again.
\end{gotcha}

\section*{A note on accuracy}

This manual was assembled by reading the source directly, and each chapter was
cross-checked against the code it describes.  Discrepancies discovered during
that process --- places where the code and its documentation or intent did not
line up --- are not silently papered over; they are recorded in
Appendix~D (Known Issues and Open Questions) for follow-up.  If you find the code
doing something this manual does not describe, Appendix~D is the first place to
look.
